\documentclass{article}
\usepackage{amsmath}
\usepackage{mathtools}
\usepackage{gensymb}
\usepackage[a4paper,inner=1.5cm,outer=1.5cm,top=2cm,bottom=0.5cm]{geometry} 
\usepackage{xcolor}
\usepackage{tikz}
\usepackage{multicol}
\usepackage{pgfplots}
\usetikzlibrary{intersections}
\usetikzlibrary{intersections,calc,angles,quotes}
\usetikzlibrary{calc,angles,positioning,intersections,quotes,decorations.markings}
\usepackage{tkz-euclide}
\usetikzlibrary{backgrounds}
\usetikzlibrary{calc,through}
\usetikzlibrary{angles}
\usetikzlibrary{fadings}
\usetikzlibrary{shapes.geometric}
\usetikzlibrary{shapes.symbols}
\usepackage{draftwatermark}
\usepackage{mathptmx}

\SetWatermarkText{\textcolor{black!30}{Mathema Shukur}}
\SetWatermarkFontSize{2 cm}
\usepackage[utf8]{inputenc}
\usepackage{fontspec}

\setmainfont{[Kalpurush.ttf]}
\newfontface{\en}{[Arial.ttf]} %%this is optional, if you want to use a secondary font. Any english font is supported
\newlength\Radius
\setlength\Radius{4cm}
\begin{document} 
	\Large
	\textcolor{red}{Welcome To} 
	\\
	\begin{tikzpicture}
		\tikz \node [fill=blue!20,star,star points=6,draw] {Mathema Shukur };
	\end{tikzpicture}
	\\
	যাদের জন্যে প্রযোজ্যঃ  	\textcolor{magenta}{একাদশ ও দ্বাদশ শ্রেণীর শিক্ষার্থী} \\
	বিষয়ঃ \textcolor{magenta}{উচ্চতর গণিত ১ম পত্র} \\
	অধ্যায়ঃ \textcolor{magenta}{৩-সরলরেখা}\\ 
	Subtopicঃ  \textcolor{magenta}{  পোলার স্থানাঙ্ক $(r,\theta)$ নির্ণয় করা }\\
	\\
	$r=\sqrt{x^2+y^2}$\\
	\\
	$\tan \theta =\frac{y}{x}$\\
	\\
	বিভিন্ন চতুর্ভাগে $\theta $ এর মান নির্ণয় \\
	\\ 
	\begin{tabular}{|c|c|c|}
		\hline
		১ম চতুর্ভাগে & $p(x,y)$ বিন্দুর জন্য & $\theta= \tan^{-1}|\frac{y}{x}|$\\
		\hline 
		২য় চতুর্ভাগে & $p(-x,y)$ বিন্দুর জন্য & $\theta= \pi - \tan^{-1}|\frac{y}{x}|$\\
		\hline
		৩য় চতুর্ভাগে & $p(-x,-y)$ বিন্দুর জন্য & $\theta= \pi+\tan^{-1}|\frac{y}{x}|,\quad \mbox{OR}\quad \theta= -\pi+\tan^{-1}|\frac{y}{x}| $\\
		\hline
		৪র্থ  চতুর্ভাগে & $p(x,-y)$ বিন্দুর জন্য & $\theta= 2\pi-\tan^{-1}|\frac{y}{x}|,\quad \mbox{OR}\quad \theta= -\tan^{-1}|\frac{y}{x}| $\\
		\hline
	\end{tabular}
	\\
	ঢাকা বোর্ড-২০২১\\
(১)	$(\sqrt{2},-\sqrt{2})$ বিন্দুর পোলার স্থানাঙ্ক নির্ণয় কর\\
	\\
	$x=\sqrt{2}$,\qquad  $y=-\sqrt{2}$\\
	\\
	$r=\sqrt{x^2+y^2}=\sqrt{(\sqrt{2})^2+(-\sqrt{2})^2}=\sqrt{2+2}=2$\\
	\\
	বিন্দুটি  ৪র্থ চতুর্ভাগ অবস্থিত\\
	\\
	$\theta= 2\pi-\tan^{-1}|\frac{y}{x}|,\quad \mbox{OR}\quad \theta= -\tan^{-1}|\frac{y}{x}| $\\
	\\ 
	$\theta= 2\pi-\tan^{-1}|\frac{-\sqrt{2}}{\sqrt{2}}|,\quad \mbox{OR}\quad \theta= -\tan^{-1}|\frac{-\sqrt{2}}{\sqrt{2}}| $\\
	\\
	$\theta= 2\pi-\tan^{-1}(1),\quad \mbox{OR}\quad \theta= -\tan^{-1}(1) $\\
	\\ 
	$\theta= 2\pi-\frac{\pi}{4},\quad \mbox{OR}\quad \theta= -\frac{\pi}{4} $\\
	\\
		$\theta= \frac{7\pi}{4},\quad \mbox{OR}\quad \theta= -\frac{\pi}{4} $\\
		\\
		$(\sqrt{2},-\sqrt{2})$ বিন্দুর পোলার স্থানাঙ্ক $(2,\frac{7\pi}{4})$ অথবা  $(2,-\frac{\pi}{4})$\\
		\\ 
	ঢাকা বোর্ড-২০২১, চট্রগ্রাম বোর্ড-২০২১, দিনাজপুর বোর্ড-২০২১ \\
(২)	$(-\sqrt{3},-1)$ বিন্দুর পোলার স্থানাঙ্ক নির্ণয় কর\\
	\\
		$x=-\sqrt{3}$,\quad $y=-1$\\
	\\
	$r=\sqrt{x^2+y^2}=\sqrt{(-\sqrt{3})^2+(-1)^2}=\sqrt{3+1}=\sqrt{4}=2$\\
	\\
	বিন্দুটি ৩য় চতুর্ভাগে অবস্থিত \\
	\\
	$\theta= \pi+\tan^{-1}|\frac{y}{x}|,\quad \mbox{OR}\quad \theta= -\pi+\tan^{-1}|\frac{y}{x}| $\\
	\\
	$\theta= \pi+\tan^{-1}|\frac{-1}{-\sqrt{3}}|,\quad \mbox{OR}\quad \theta= -\pi+\tan^{-1}|\frac{-1}{-\sqrt{3}}| $\\
	\\
	$\theta= \pi+\tan^{-1}\frac{1}{\sqrt{3}},\quad \mbox{OR}\quad \theta= -\pi+\tan^{-1}\frac{1}{\sqrt{3}} $\\
	\\
	$\theta= \pi+\frac{\pi}{6},\quad \mbox{OR}\quad \theta= -\pi+\frac{\pi}{6} $\\
	\\
	$\theta=\frac{7\pi}{6},\quad \mbox{OR}\quad \theta= \frac{-5\pi}{6} $\\
	\\
$(-\sqrt{3},-1)$	বিন্দুটির পোলার স্থানাঙ্ক $(r,\theta)=(2,\frac{7\pi}{6}),\qquad \mbox{OR}\quad (2,-\frac{5\pi}{6})$\\
	\\
	যশোর বোর্ড-২০২১\\
(৩)	$(-3,3)$ বিন্দুর পোলার স্থানাঙ্ক নির্ণয় কর\\
	\\
	$x=-3$,\qquad  $y=3$\\
	\\
	$r=\sqrt{x^2+y^2}=\sqrt{(-3)^2+(3)^2}=\sqrt{9+9}=3\sqrt{2}$\\
	\\
	বিন্দুটি ২য় চতুর্ভাগ অবস্থিত\\
	\\ 
	$\theta= \pi - \tan^{-1}|\frac{y}{x}|$\\
	\\
	$\theta= \pi - \tan^{-1}|\frac{3}{-3}|$\\
	\\
	$\theta= \pi - \tan^{-1}1$\\
	\\
	$\theta= \pi - \frac{\pi}{4}$\\
	\\
	$\theta=\frac{3\pi}{4}$\\
	\\
	$(-3,3)$	বিন্দুটির পোলার স্থানাঙ্ক $(r,\theta)=(3\sqrt{2},\frac{3\pi}{4})$\\
	\\
	যশোর বোর্ড-২০২১\\
	(৪) $(-3,-3)$ বিন্দুর পোলার স্থানাঙ্ক নির্ণয় কর\\
	$\textcolor{red}{(r,\theta)=(3\sqrt{2},\frac{5\pi}{4})}$ অথবা ,\qquad $\textcolor{red}{(r,\theta)=(3\sqrt{2},-\frac{3\pi}{4})}$
	\\
	সিলেট বোর্ড-২০২১\\
(৫)	$(-\frac{3\sqrt{2}}{2},\frac{3\sqrt{2}}{2})$ বিন্দুর পোলার স্থানাঙ্ক নির্ণয় কর\\
	$\textcolor{red}{(r,\theta)=(3,\frac{3\pi}{4})}$\\
	\\
		(৬) রাজশাহী বোর্ড-২০১৯\\
	$x+3y+3=0$ রেখাটি y- অক্ষকে যে বিন্দুতে ছেদ করে তাঁর পোলার স্থানাঙ্ক নির্ণয় কর। \\
	ছেদবিন্দু $(0,-1)$\\ 
	$\textcolor{red}{(r,\theta)=(1,\frac{3\pi}{2})}$ , অথবা	$\textcolor{red}{(r,\theta)=(1,-\frac{\pi}{2})}$\\ 
	\\
	(৭) যশোর বোর্ড-২০১৯; সকল বোর্ড-২০১৮;ঢাকা বোর্ড-২০১৭;দিনাজপুর বোর্ড-২০১৭\\
	$(-1,\sqrt{3})$ এর পোলার স্থানাঙ্ক নির্ণয় কর। \\
	$\textcolor{red}{(r,\theta)=(2,\frac{2\pi}{3})}$\\
	\\
	(৮) রাজশাহী বোর্ড-২০২১, বরিশাল বোর্ড-২০১৯\\
	$(-1,-1)$ এর পোলার স্থানাঙ্ক নির্ণয় কর।\\
	$\textcolor{red}{(r,\theta)=(\sqrt{2},\frac{5\pi}{4})}$ অথবা ,\qquad $\textcolor{red}{(r,\theta)=(\sqrt{2},-\frac{3\pi}{4})}$\\
	\\
	(৯) ঢাকা বোর্ড-২০১৯\\
	$(3,\sqrt{3})$ এর পোলার স্থানাঙ্ক নির্ণয় কর। \\
	$\textcolor{red}{(r,\theta)=(2\sqrt{3},\frac{\pi}{6})}$\\
	\\
	(১০) কুমিল্লা বোর্ড -২০২১, সকল বোর্ড-২০১৮\\
	$(-4,-4)$ এর পোলার স্থানাঙ্ক নির্ণয় কর। \\
	$\textcolor{red}{(r,\theta)=(4\sqrt{2},\frac{5\pi}{4})}$ অথবা ,\qquad $\textcolor{red}{(r,\theta)=(4\sqrt{2},-\frac{3\pi}{4})}$\\
	\\
	(১১) সিলেট বোর্ড-২০১৭\\
	$(1,-1)$ এর পোলার স্থানাঙ্ক নির্ণয় কর। \\ 
	$\textcolor{red}{(r,\theta)=(\sqrt{2},\frac{7\pi}{4})}$ অথবা  $\textcolor{red}{(r,\theta)=(\sqrt{2},-\frac{\pi}{4})}$\\
	\\  
	(১২) সিলেট বোর্ড-২০১৭\\
	$(-2,-\sqrt{2})$ এর পোলার স্থানাঙ্ক নির্ণয় কর। \\
	$\textcolor{red}{(r,\theta)=(\sqrt{6},\tan^{-1}\frac{\sqrt{2}}{2})}$ অথবা $\textcolor{red}{(r,\theta)=(\sqrt{6},-\pi+\tan^{-1}\frac{\sqrt{2}}{2})}$\\
	\\
	(১৩) কুমিল্লা বোর্ড-২০১৭\\
	$(-\sqrt{2},-\sqrt{2})$ এর পোলার স্থানাঙ্ক নির্ণয় কর। \\ 
	$\textcolor{red}{(r,\theta)=(2,\frac{5\pi}{4})}$ অথবা ,\qquad $\textcolor{red}{(r,\theta)=(2,-\frac{3\pi}{4})}$  
	\end{document}